\section*{Chapter \ref{ch:ll:the-feed-handler} --- The Feed Handler}

\begin{solution}{exo:ll:the-feed-handler:1}
Line A's packet holds 1\,003--1\,005 and opens a gap at 1\,000; line B's brings 1\,000--1\,002, which are published at once,
followed by the held 1\,003--1\,005: six events in sequence order. The book was stale from the first packet's arrival to the
second's, \qty{40}{\micro\second}, and the gap counts as filled by the other line.
\end{solution}

\begin{solution}{exo:ll:the-feed-handler:2}
512~KiB (the kernel doubles the value asked). At 960 bytes per 64-byte datagram it holds 546 of them, as measured.
\end{solution}

\begin{solution}{exo:ll:the-feed-handler:3}
The next snapshot comes on average half an interval later and at worst a full interval: $250 + 5 = \qty{255}{\milli\second}$ on
average, $500 + 5 = \qty{505}{\milli\second}$ at worst.
\end{solution}

\begin{solution}{exo:ll:the-feed-handler:4}
Most gaps are messages that the other line delivers a few microseconds later or that a switch reordered: in the chapter's
minute, 4\,224 of 4\,232. Asking at once would load the server and the network with requests for messages already on their
way. A shorter wait sends more useless requests; a longer one leaves the book stale longer whenever a message really is lost
on both lines.
\end{solution}

\begin{solution}{exo:ll:the-feed-handler:5}
$(2 - 1.2) \times 10^6 \times 0.03 = 24\,000$ datagrams queue. The kernel must grant $24\,000 \times 960 \approx 23$~MB, so
about 11.5~MB must be asked, which needs \texttt{rmem\_max} raised well above the laptop's 4~MiB.
\end{solution}

\begin{solution}{exo:ll:the-feed-handler:6}
The snapshot replaces the lost messages by the state they led to: the book after the snapshot equals the clean run's book at
that sequence number, and every event after it is the clean run's event. The build checks both: the events after the
snapshot's sequence number equal the clean run's, and the final book built from each stream is the same.
\end{solution}

\begin{solution}{exo:ll:the-feed-handler:7}
Collect the holes between the next expected number and the largest held one (the gaps in the held map's keys) and send them
as one request with a list of ranges; the server answers all of them in one round trip. The \qty{2.7}{\milli\second} gap of the
minute, which needed several rounds, should then close in one: the timeout plus one round trip, \qty{0.7}{\milli\second}.
\end{solution}

\begin{solution}{exo:ll:the-feed-handler:8}
Most gaps are filled by the other line within microseconds or by a retransmission within a millisecond; waiting for a snapshot
turns each of them into a stale book for up to a snapshot interval, a thousand times longer: in the chapter's minute, more
than four thousand gaps would each have waited for a snapshot. Consistency comes from holding messages until the gap is closed, not from discarding them.
\end{solution}

\begin{solution}{pb:ll:the-feed-handler:1}
\begin{enumerate}
\item 500\,000 datagrams a second; the burst asks for three times that, a utilisation of 3.
\item $(1.5 - 0.5) \times 10^6 \times 0.02 = 20\,000$.
\item The queue shrinks at $500\,000 - 200\,000 = 300\,000$ a second: about \qty{67}{\milli\second}.
\item It waits behind 20\,000 others served at \qty{2}{\micro\second}: \qty{40}{\milli\second}.
\item $20\,000 \times 960 = 19.2$~MB granted.
\item 9.6~MB, more than \texttt{rmem\_max}: the kernel caps it at 4~MiB (8~MiB granted) unless the administrator raises the
limit or the process has the privilege to force it.
\item The 8~MiB buffer holds 8\,738 datagrams: about 11\,262 are lost.
\item Memory the handler controls: a reader thread that only moves datagrams from the socket to a large ring (the work is done
by another thread), or the network card's own queues read directly (\omterm{def:ll:linux-tuning:polling}{kernel bypass}, chapter~13); or a faster handler.
\item Yes if the server allows it: 11\,262 messages are within its 100\,000-message window, in twelve requests of at most 1\,000.
\item Half a snapshot interval on average, \qty{0.5}{\second}, and a whole one at worst, \qty{1}{\second}, plus loading.
\item At 200\,000 a second after the burst, about 100\,000 on average and 200\,000 at worst.
\item Stop quoting on the instrument and cancel resting orders that depend on the book, until the flag is lowered.
\item \textbf{Named result.} Absorbing the burst takes 20\,000 datagrams of queue, 19.2~MB of receive buffer granted (9.6~MB
asked, more than this machine's 4~MiB limit); when it fails and only a snapshot can recover, the book is stale for
\qty{0.5}{\second} on average and \qty{1}{\second} at worst with a snapshot a second.
\item Three times faster: serving 1.5 million a second, \qty{0.67}{\micro\second} a datagram.
\item Each line has its own socket and buffer, so the memory doubles, and the handler must read twice the datagrams, half of
them duplicates discarded after a comparison.
\item At 350 million messages a second in its busiest millisecond, a feed of that size cannot be handled by one thread at tens
of nanoseconds a message: it must be split across lines and cores.
\item The distribution of arrivals per 10 milliseconds and per millisecond from captures of the busiest days, the handler's
service time under that load, and the kernel's drop counters for the sockets.
\item The receive queue's high-water mark (and any drop counted by the kernel).
\item A shorter interval costs the venue bandwidth on a channel most subscribers never read; the firm pays for a longer one in
staleness after every unrecoverable gap.
\item It queues what it can, loses the rest, says so with its stale flag, and recovers from retransmission or the next
snapshot.
\end{enumerate}
\end{solution}

\begin{solution}{iq:ll:the-feed-handler:1}
It receives the redundant lines (from the network card through the kernel or directly), decodes the packets, arbitrates by
sequence number, detects and recovers gaps (other line, retransmission, snapshot), normalises each message into the firm's
record, publishes it on shared-memory rings with a staleness flag, and keeps counters for monitoring.
\iqlookfor{the whole chain, with recovery and a staleness signal.}
\end{solution}

\begin{solution}{iq:ll:the-feed-handler:2}
Hold the later messages and raise the stale flag; wait briefly for line B; if the gap is not filled, ask the retransmission
server for the range and apply the answer in order; if the server cannot help, wait for the instrument's next snapshot, load
it, discard held messages it already reflects, apply the rest, and lower the flag.
\iqlookfor{the three recovery paths in order, and the flag.}
\end{solution}

\begin{solution}{iq:ll:the-feed-handler:3}
From the largest burst expected: the backlog (arrival rate minus service rate, times the burst's duration) times what the kernel
charges per datagram, much more than its payload. The kernel doubles the value set and caps it at \texttt{rmem\_max};
check the granted value with \texttt{getsockopt} and watch drop counters.
\iqlookfor{sizing from data, the doubling, the cap and the per-datagram charge.}
\end{solution}

\begin{solution}{iq:ll:the-feed-handler:4}
The handler publishes its \omterm{def:ll:the-feed-handler:stale}{stale-book flag} with the events (and the sequence numbers let a consumer see a snapshot reset). A
consumer that keeps its own book checks the flag before using it, and checks invariants such as an uncrossed book.
\iqlookfor{an explicit flag propagated with the data.}
\end{solution}

\begin{solution}{iq:ll:the-feed-handler:5}
Downstream code (books, strategies, logs) is written once for every venue, records have a fixed size for rings and logs, and
sequence numbers travel with them. The cost is a copy per message and fields that some venues do not fill; a venue's
specifics must be carried in the record or they are lost.
\iqlookfor{one downstream interface against a copy and lost detail.}
\end{solution}

\begin{solution}{iq:ll:the-feed-handler:6}
Replay recorded lines with injected losses, reorderings, outages on one and both lines, and bursts, against a retransmission
server and a \omterm{def:ll:protocols-ii-binary-exchange-protocols:channels}{snapshot channel} that can be made to fail; compare the output with a clean run (identical events after
retransmission; identical book and later events after a snapshot), in every language the handler is written in, as the build
does.
\iqlookfor{deterministic replay with injected faults and a clean reference.}
\end{solution}
